\scanpage{22}
\deleted{\illegible{1} $A\in\mathcal C$. La construction}

Dans le cas où $\mathcal C$ est une catégorie \uline{abélienne} ayant déjà suffisamment d'injectifs, la question se pose de savoir quand les foncteurs dérivés droits \inserted{usuels} d'un foncteur $\Gamma:\mathcal C\unsure{-|\to}\mathcal D$ coïncident avec les foncteurs dérivés usuels. \deleted{Voici des} \inserted{Voici des} conditions suffisantes raisonnables. On veut tout au moins que $R_S^*\Gamma$ soit un foncteur cohomologique,

\deleted{Supposons que $\varphi$ soit injectif, i.e. $A\to TS(A)$ injectif pour tout $A\in\mathcal C$. \unsure{Supposons} $\Gamma$ exact \unsure{à gauche}. \illegible{2} $S$ et $T$ exacts, donc $TS$ exact, donc que la résolution canonique $C(A)$\unsure{?} La $S$-résolution canonique de $A$ est une résolution au sens usuel, d'où résulte que $R^0_S\Gamma=\Gamma=R^0\Gamma$. Il \unsure{reste|vaut} donc à prouver l'effaçabilité des $R_S^p\Gamma$}

i.e. l'opérateur cobord $\partial$ est défini pour \uline{\unsure{ttes}} les suites exactes en $A'-A-A''-0$ (et ``non seulement'' les suites $S$-exactes). C'est \unsure{le|le le|donc le} cas si $\Gamma(C^*(A))$ est un foncteur exact en $A$, pour ceci nous supposons.

% margin note written sideways
\margin{Ces conditions d'ailleurs suffisent, car \deleted{\unsure{alors}} $(R_S^*\Gamma)$ sera aussi \unsure{coeffaçable|effaçable}, car il s'annule \deleted{\unsure{évidemment}} sur les objets de la forme $T(X)$, qui sont ``assez nombreux'' grâce à (i). \deleted{soit un foncteur \uline{exact} par r. à $A$.}}

(i) $A\to C^0(A)$ injectif

(ii) $S$ et $T$ [donc $C^0$, donc \inserted{aussi} $C^*$ \inserted{compte tenu de (i)}] foncteurs exacts

(iii) Si $0\to A'-A-A''-0$ est une suite exacte de $\mathcal C$, alors la suite $0-\Gamma C^0(A')\to\Gamma C^0(A)-\Gamma C^0(A'')\inserted{-0}$ est exacte [Il suffit que le foncteur \uuline{$\Gamma T$ soit exact}]

\deleted{Enfin, pour} \inserted{Alors} $R^*_S\Gamma$ est bien un foncteur cohomologique. Pour qu'en degré $0$, il coïncide avec $R^0\Gamma$, on va supposer que

\unsure{(iv)} \deleted{\unsure{les} \illegible{1} $A\to C^0(A)$ injectifs} $\Gamma$ exact à gauche.

% archive sheet number: 21
